% ===================================================================
% Section 2: Popis systému + Identifikácia 
% ===================================================================

\section{Identifikácia a popis tepelného systému TS05}
\label{sec:system_ident_ts05}

V tejto kapitole sa budeme zaoberať laboratórnym modelom, ktorý bol poskytnutý pre potreby tejto diplomovej práce – tepelným systémom TS05.  
Cieľom kapitoly je podrobne opísať túto experimentálnu zostavu, vykonať jej identifikáciu a navrhnúť viaceré riešenia riadenia, a to ako klasické, tak aj pokročilejšie metódy.

Súčasťou práce je tiež implementácia vybraných riadiacich algoritmov na rôznych hardvérových platformách.  
Táto časť má umožniť porovnanie prístupov k implementácii riadenia v simulačnom prostredí MATLAB/Simulink a na vybranom mikrokontroléri, čím sa vytvorí priestor pre analýzu rozdielov v správaní systému a v samotnej realizácii algoritmov.

% -------------------------------------------------------------------
% Subsection 2.1 - Tepelný systém TS05 
% -------------------------------------------------------------------

\subsection{Popis tepelného systému}

Táto kapitola sa zaoberá popisom laboratórneho modelu TS05 a jeho základnými statickými vlastnosťami.  
Model predstavuje fyzický systém určený na výučbu a experimentálne overovanie princípov riadenia v oblasti automatizácie a mechatroniky.

% -------------------------------------------------------------------
% Subsubsection 2.1.1 - Popis systému
% -------------------------------------------------------------------

\subsubsection{Popis laboratórneho fyzikálneho modelu}

Laboratórny model TS05 predstavuje MIMO systém (Multiple Input – Multiple Output) so dvomi vstupmi a dvomi výstupmi, ktorého zapojenie je uvedené v tabuľke~\ref{tab:io_ts05}.  
Vstupné veličiny tvoria napätia riadiace ventilátor a vykurovaciu špirálu (v rozsahu 0–10~V), zatiaľ čo výstupy sú merané pomocou Senzora~1 a Senzora~2 (tiež v rozsahu 0–10~V).

\begin{table}[h!]
    \centering
    \caption{Vstupy a výstupy systému TS05}
    \label{tab:io_ts05}
    \begin{tabular}{|c|c|c|}
        \hline
        \textbf{Typ signálu} & \textbf{Označenie} & \textbf{Rozsah} \\ \hline
        Vstup & Ventilátor & 0–10~V \\ \hline
        Vstup & Výkon špirály & 0–10~V \\ \hline
        Výstup & Senzor~1 (teplota 1) & 0–10~V \\ \hline
        Výstup & Senzor~2 (teplota 2) & 0–10~V \\ \hline
    \end{tabular}
\end{table}

Vizuálny vzhľad systému je znázornený na obrázku~\ref{fig:TS_sch3_foto}.

\begin{center}
\vbox{
\makebox[\textwidth][c]{%
    \includegraphics[width=1.1\textwidth]{TS_sch3_foto.pdf}
}
\figcaption{Laboratórny model TS05 – reálna zostava~\cite{kut015}}
\label{fig:TS_sch3_foto}
}
\end{center}

Systém komunikuje s počítačom prostredníctvom meracej a komunikačnej karty \textit{Humusoft},  
ktorá je určená na prepojenie reálnych systémov so simulačným prostredím MATLAB/Simulink.  
Táto karta umožňuje obojsmerný prenos analógových signálov – prijíma merané veličiny zo systému a zároveň vysiela riadiace signály na jeho vstupy.

Samotná komunikácia a spracovanie signálov prebieha v prostredí Simulink,  
v rámci vopred pripravenej schémy (obr.~\ref{fig:TS_simulink_scheme}),  
ktorá bola vytvorená ako súčasť projektu KUT – Krátke výučbové materiály.  
Táto schéma slúži aj ako podklad pre predmet \textit{Modelovanie a riadenie systémov}~\cite{kut015}.

\begin{center}
\vbox{
\makebox[\textwidth][c]{%
    \includegraphics[width=1.3\textwidth, trim=100 200 100 200, clip]{TS_simulink_scheme.pdf}
}
\figcaption{Simulinková schéma pre komunikáciu so systémom TS05}
\label{fig:TS_simulink_scheme}
}
\end{center}

% -------------------------------------------------------------------
% Subsubsection 2.1.2 - Statické vlastnosti systému
% -------------------------------------------------------------------

\subsubsection{Statické vlastnosti tepelného systému}
\label{sec:staticke_vlastnosti}

Na lepšie pochopenie správania systému boli experimentálne namerané jeho statické charakteristiky.  
Počas meraní boli skúmané tri pracovné body ventilátora – 3~V, 6~V a 9~V – pričom v každom prípade bola hodnota napätia ventilátora udržiavaná konštantná.  
Napätie na vykurovacej špirále sa menilo lineárne od 0~V do 10~V počas celého merania, keďže špirála má výraznejší vplyv na teplotu ako ventilátor.  
Priebeh signálu špirály počas experimentov je zobrazený na obrázku~\ref{fig:ts_prevodova_spirala}.

\begin{center}
\vbox{
\makebox[\textwidth][c]{%
    \includegraphics[width=1.5\textwidth]{ts_prevodova_spirala.png}
}
\figcaption{Priebeh napätia na špirále počas všetkých troch experimentov}
\label{fig:ts_prevodova_spirala}
}
\end{center}

Výsledkom meraní sú šesť grafov, ktoré zobrazujú odozvu systému na zmenu vstupných veličín.  
Každý graf obsahuje výstupy zo Senzora~1 a Senzora~2 ako funkcie času pre jednotlivé pracovné body ventilátora.

\begin{center}
\vbox{
\makebox[\textwidth][c]{%
    \includegraphics[width=1.5\textwidth]{ts_prevodova_sensors3.png}
}
\figcaption{Odozva systému TS05 pre pracovný bod ventilátora 3~V}
\label{fig:ts_prevodova_sensors3}
}
\end{center}

\begin{center}
\vbox{
\makebox[\textwidth][c]{%
    \includegraphics[width=1.5\textwidth]{ts_prevodova_sensors6.png}
}
\figcaption{Odozva systému TS05 pre pracovný bod ventilátora 6~V}
\label{fig:ts_prevodova_sensors6}
}
\end{center}

\begin{center}
\vbox{
\makebox[\textwidth][c]{%
    \includegraphics[width=1.5\textwidth]{ts_prevodova_sensors9.png}
}
\figcaption{Odozva systému TS05 pre pracovný bod ventilátora 9~V}
\label{fig:ts_prevodova_sensors9}
}
\end{center}

Z tých istých meraní boli zostavené aj prevodové (statické) charakteristiky,  
ktoré opisujú závislosť výstupných veličín od vstupného napätia špirály pre jednotlivé pracovné body ventilátora.

\begin{center}
\vbox{
\makebox[\textwidth][c]{%
    \includegraphics[width=1.5\textwidth]{ts_static_sensors3.png}
}
\figcaption{Prevodová charakteristika systému TS05 pre pracovný bod ventilátora 3~V}
\label{fig:ts_static_sensors3}
}
\end{center}

\begin{center}
\vbox{
\makebox[\textwidth][c]{%
    \includegraphics[width=1.5\textwidth]{ts_static_sensors6.png}
}
\figcaption{Prevodová charakteristika systému TS05 pre pracovný bod ventilátora 6~V}
\label{fig:ts_static_sensors6}
}
\end{center}

\begin{center}
\vbox{
\makebox[\textwidth][c]{%
    \includegraphics[width=1.5\textwidth]{ts_static_sensors9.png}
}
\figcaption{Prevodová charakteristika systému TS05 pre pracovný bod ventilátora 9~V}
\label{fig:ts_static_sensors9}
}
\end{center}

% -------------------------------------------------------------------
% Subsection 2.2 - Identifikácia systému TS05 
% -------------------------------------------------------------------

\subsection{Identifikácia tepelného systému}
\label{sec:identifikacia}

V tejto časti sa pokúsime identifikovať tepelný systém TS05 rôznymi metódami a porovnať ich výsledky.  
Cieľom je určiť, ktoré prístupy poskytujú presnejšie modely a na aké faktory je potrebné dávať pozor počas procesu identifikácie.

% -------------------------------------------------------------------
% Subsubsection 2.2.1 - Meranie údajov a ich spracovanie
% -------------------------------------------------------------------

\subsubsection{Meranie údajov a ich spracovanie}

Na účely identifikácie systému akoukoľvek \textit{off-line} metódou je najprv potrebné získať reakciu systému na jednotkové skoky v okolí pracovného bodu.  
V našom prípade bol pracovný bod zvolený ako vstupné napätia špirály~4~V a ventilátora~6~V.  

Na obrázku~\ref{fig:ts_pb_vent6_spir4_ident_raw} je zobrazený vstupný signál špirály a odpoveď systému na výstupe zo senzora~1.

\begin{center}
\vbox{
\makebox[\textwidth][c]{%
    \includegraphics[width=1.5\textwidth]{ts_pb_vent6_spir4_ident_raw.png}
}
\figcaption{Vstupný signál špirály a výstup senzora~1 pri skokovej zmene}
\label{fig:ts_pb_vent6_spir4_ident_raw}
}
\end{center}

Z grafu možno pozorovať zreteľný lineárny trend, teda postupné zvyšovanie výstupnej teploty počas celého merania.  
Tento jav je spôsobený pomalým nahrievaním špirály a tiež prirodzenými zmenami teploty v miestnosti, ktoré neboli v rámci experimentu kontrolované.  
Teplota prostredia sa však počas meraní pohybovala v rozsahu bežnej izbovej teploty, takže systém bol v stabilných podmienkach.

Aby bolo možné získať z nameraných údajov vhodné parametre pre identifikáciu, bolo potrebné odstrániť tento trend aplikovaním odstránenia lineárneho trendu.  
Metodika výpočtu je založená na lineárnej regresii, kde sa určuje priamka trendu v tvare

\begin{equation}
    \hat{y} = \hat{\beta}_0 + \hat{\beta}_1 x
    \label{eq:linear_trend}
\end{equation}

pričom koeficienty sa počítajú podľa vzťahov~\cite{montgomery}:

\begin{equation}
    \hat{\beta}_1 = \frac{n\sum x_i y_i - \sum x_i \sum y_i}{n\sum x_i^2 - (\sum x_i)^2},
    \qquad
    \hat{\beta}_0 = \bar{y} - \hat{\beta}_1 \bar{x}
    \label{eq:beta_coefficients}
\end{equation}

Takto určený trend je následne odčítaný od pôvodného signálu:

\begin{equation}
    y_{\text{detr}} = y - (\hat{\beta}_0 + \hat{\beta}_1 x)
    \label{eq:detrended_signal}
\end{equation}

Uvedený postup bol implementovaný v prostredí \textit{MATLAB} pomocou vlastnej funkcie:

\begin{lstlisting}[language=Matlab, 
caption={MATLAB implementácia detrendácie signálu pomocou lineárnej regresie},
label={lst:my_detrend}]
function y_detr = my_detrend(y, x)
    y = y(:);
    x = x(:);
    n = length(y);

    Sx  = sum(x);
    Sy  = sum(y);
    Sxx = sum(x.^2);
    Sxy = sum(x .* y);

    b1 = (n*Sxy - Sx*Sy) / (n*Sxx - Sx^2);  
    b0 = mean(y) - b1 * mean(x);            

    y_detr = y - (b0 + b1*x);
end
\end{lstlisting}

Po odstránení lineárneho trendu bol signál zároveň normalizovaný pre potreby ďalšej identifikácie.  
Výsledné priebehy normalizovaných a detrendovaných údajov sú zobrazené na obrázku~\ref{fig:ts_pb_vent6_spir4_ident_raw_normalized}.

\begin{center}
\vbox{
\makebox[\textwidth][c]{%
    \includegraphics[width=1.5\textwidth]{ts_pb_vent6_spir4_ident_raw_normalized.png}
}
\figcaption{Normalizované a detrendované priebehy vstupného a výstupného signálu}
\label{fig:ts_pb_vent6_spir4_ident_raw_normalized}
}
\end{center}

Z grafov možno vidieť, že jednotlivé skokové odozvy sa mierne líšia svojím priebehom,  
čo poukazuje na citlivosť systému TS05 na vonkajšie podmienky a fluktuácie prostredia.  
Preto boli vykonané tri nezávislé experimenty s rovnakým vstupným signálom,  
ktorých priebehy sú zobrazené na obrázku~\ref{fig:ts_pb_vent6_spir4_ident_raw_normalized}.  
Po detrendácii a normalizácii boli tieto dáta spriemerované bod po bode,  
aby sa eliminoval vplyv náhodných rušení a získala reprezentatívna odozva systému.  

Výsledný stredný priebeh, znázornený na obrázku~\ref{fig:ts_pb_vent6_spir4_ident_avg},  
predstavuje typickú reakciu systému TS05 v okolí pracovného bodu 4~V s rozsahom~$\pm 0{,}5$~V  
a slúži ako podklad pre ďalšiu identifikáciu modelu.

\begin{center}
\vbox{
\makebox[\textwidth][c]{%
    \includegraphics[width=1.5\textwidth]{ts_pb_vent6_spir4_ident_avg.png}
}
\figcaption{Priemerná odozva systému TS05 v pracovnom bode 4~V}
\label{fig:ts_pb_vent6_spir4_ident_avg}
}
\end{center}

% -------------------------------------------------------------------
% Subsubsection 2.2.2 - Identifikácia pomocou metódy najmenších štvorcov
% -------------------------------------------------------------------

\subsubsection{Identifikácia systému pomocou metódy najmenších štvorcov}

V tejto časti bola uskutočnená identifikácia systému TS05 na základe priemernej odozvy získanej z troch experimentálnych meraní popísaných v predchádzajúcej kapitole.  
Tento priemerný priebeh reprezentuje typickú reakciu systému v okolí pracovného bodu 4~V a slúži ako základ pre odhad parametrov modelu.  
Cieľom bolo získať lineárny model druhého rádu, ktorý dostatočne presne opisuje dynamické správanie systému v tejto oblasti.

Identifikácia bola realizovaná pomocou metódy najmenších štvorcov~\cite{fikar},  
ktorá umožňuje odhad parametrov diskrétneho modelu systému na základe meraných vstupov a výstupov.  
Pre odhad parametrov bola vytvorená regresná matica $H(k,:)$ a vektor výstupov $Y$, pričom výsledný odhad parametrov $\hat{\theta}$ bol určený podľa rovníc:

\begin{equation}
    H(k,:) = [-y(k-1), -y(k-2), u(k-1), u(k-2)]
    \label{eq:phi_matrix}
\end{equation}

\begin{equation}
    \hat{\theta} = (H^\mathrm{T}H)^{-1}H^\mathrm{T}Y
    \label{eq:theta_estimation}
\end{equation}

Na základe týchto odhadov boli vypočítané koeficienty diskrétneho modelu druhého rádu,  
ktorý následne predstavuje prenosovú funkciu $G_d(z)$.  
Pre účely simulácie a ďalšieho návrhu riadenia bol tento model prevedený na spojitú formu $G(s)$ pomocou Tustinovej aproximácie.

V prostredí MATLAB bola implementovaná funkcia, ktorá realizuje uvedené výpočty a vytvára identifikovaný model systému.  
Implementácia tejto funkcie je uvedená nižšie:

\begin{lstlisting}[language=Matlab, 
caption={MATLAB implementácia identifikácie metódou najmenších štvorcov},
label={lst:ls_identify_basic}]
function G = ls_identify_basic(y, u, Ts)
    
    y = y(:);
    u = u(:);
    N = length(y);
    na = 2; nb = 2;
    
    H = zeros(N-2, na+nb);
    for k = 3:N
        H(k-2,:) = [-y(k-1), -y(k-2), u(k-1), u(k-2)];
    end
    Y = y(3:N);
    
    theta = pinv(H) * Y;  
    
    a1 = theta(1); 
    a2 = theta(2);
    b1 = theta(3); 
    b2 = theta(4);
    
    num = [b1 b2];
    den = [1 a1 a2];
    Gd = tf(num, den, Ts);
    
    G = d2c(Gd, 'tustin');

end
\end{lstlisting}

% -------------------------------------------------------------------
% Subsubsection 2.2.3 - Identifikácia pomocou optimalizácie 
% -------------------------------------------------------------------

\subsubsection{Identifikácia systému pomocou optimalizácie prenosovej funkcie}

V tejto časti bol zvolený menej akademický, skôr prakticky orientovaný prístup k identifikácii systému TS05.  
Namiesto priameho odhadu parametrov modelu metódou najmenších štvorcov bola použitá optimalizácia parametrov prenosovej funkcie tak,  
aby sa minimalizovala chyba medzi simulovanou odozvou modelu a reálnou priemernou odozvou systému, získanou z experimentov.

Ako nástroj optimalizácie bola použitá funkcia \texttt{fminsearch} z prostredia MATLAB,  
ktorá implementuje Nelderovu–Meadovu simplexovú metódu.  
Ide o bezderivačnú lokálnu optimalizačnú metódu, ktorá na základe hodnotenia účelovej funkcie v niekoľkých bodoch parametrov iteratívne posúva simplex (mnohouholník v priestore parametrov) tak, aby sa znižovala hodnota účelovej funkcie.  
V našom prípade je účelová funkcia definovaná ako koreň zo strednej kvadratickej chyby (RMSE) medzi nameranou odozvou systému a simulovaným výstupom modelu:

\begin{equation}
    J(\mathbf{p}) = 
    \sqrt{
        \frac{1}{N} 
        \sum_{k=1}^{N} 
        \bigl(y(k) - y_{\text{sim}}(k,\mathbf{p})\bigr)^2
    }
    \label{eq:rmse_function}
\end{equation}

kde $y(k)$ je nameraný výstup, $y_{\text{sim}}(k,\mathbf{p})$ je výstup simulovaný modelom s parametrami $\mathbf{p}$  
a $N$ je počet vzoriek merania.

Optimalizácia prebieha nad parametrami diskrétneho modelu druhého rádu, ktoré sú následne prevedené na spojitú prenosovú funkciu pomocou Tustinovej aproximácie.  
Implementácia použitej funkcie v prostredí MATLAB je uvedená nižšie:

\begin{lstlisting}[language=Matlab,
caption={MATLAB implementácia identifikácie pomocou optimalizácie prenosovej funkcie},
label={lst:tf_identify_fmin}]
function G = tf_identify_fmin(y, u, time, Ts)

    y = y(:);
    u = u(:);
    time = time(:);

    p0 = [0 0 0 0];

    cost_fun = @(p) rmse_cost_stable(p, u, y, time, Ts);

    opts = optimset('Display', 'off', 'TolX', 1e-8, 'TolFun', 1e-8);
    p_opt = fminsearch(cost_fun, p0, opts);

    p_opt = 50 * tanh(p_opt);
    a1 = p_opt(1); a2 = p_opt(2); b1 = p_opt(3); b2 = p_opt(4);

    Gd = tf([b1 b2], [1 a1 a2], Ts);
    G  = d2c(Gd, 'tustin');

end

function err = rmse_cost_stable(p, u, y, t, Ts)

    p = 50 * tanh(p);
    a1 = p(1); a2 = p(2); b1 = p(3); b2 = p(4);

    try

        Gd = tf([b1 b2], [1 a1 a2], Ts);

        poles = pole(Gd);
        if any(abs(poles) >= 1)
            err = 1e6;
            return;
        end

        G = d2c(Gd, 'tustin');
        y_sim = lsim(G, u, t);

        err = sqrt(mean((y - y_sim).^2));

    catch
        err = 1e6;
    end
end
\end{lstlisting}

Pri tomto prístupe je dôležité upozorniť na jeden podstatný problém.  
Ak nie sú parametre prenosovej funkcie počas optimalizácie vhodne obmedzené,  
algoritmus môže nájsť také hodnoty parametrov, ktoré síce veľmi dobre aproximujú konkrétnu meranú odozvu,  
ale vedú k nerealistickému alebo nestabilnému správaniu modelu pri iných vstupoch.  
V predchádzajúcich pokusoch viedlo práve toto k prenosovým funkciám s „nezdravými“ parametrami,  
ktoré boli použiteľné iba pre konkrétny fitovaný experiment.

Z tohto dôvodu sú v uvedenej implementácii parametre počas optimalizácie mapované pomocou transformácie

\begin{equation}
	p_i^\ast = 50 \tanh(p_i)
	\label{eq:param_transform_extended}
\end{equation}

čím sú efektívne udržiavané v konečnom rozsahu a znižuje sa riziko vzniku numericky extrémnych alebo fyzikálne nezmyselných modelov.  
Okrem toho bola do optimalizácie pridaná aj dodatočná penalizácia, ktorá výrazne zvyšuje hodnotu účelovej funkcie v prípade,  
že identifikovaná prenosová funkcia vykazuje známky nestability.  
Týmto spôsobom algoritmus uprednostňuje iba stabilné riešenia a zabraňuje vzniku modelov so „zdravotne“ nevhodnými pólmi.

% -------------------------------------------------------------------
% Subsubsection 2.2.4 - Výsledky identifikácie TS05
% -------------------------------------------------------------------

\subsubsection{Výsledky identifikácie tepelného systému}

Po identifikácii systému v pracovnom bode pomocou dvoch rôznych metód  
(metódy najmenších štvorcov a optimalizácie prenosovej funkcie)  
boli získané nasledujúce výsledky.

\begin{center}
\vbox{
\makebox[\textwidth][c]{%
    \includegraphics[width=1.5\textwidth]{ts_pb_vent6_spir4_ident_comparison.png}
}
\figcaption{Porovnanie skokovej odozvy reálneho systému TS05 a identifikovaných modelov}
\label{fig:ts_pb_vent6_spir4_ident_comparison}
}
\end{center}

Porovnanie priemernej skokovej odozvy reálneho systému a odoziev oboch identifikovaných modelov  
je zobrazené na obrázku~\ref{fig:ts_pb_vent6_spir4_ident_comparison}.

Identifikovaný model získaný metódou najmenších štvorcov má tvar

\begin{equation}
    G_{s,\mathrm{LS}}(s) = 
    \frac{0.00171\,s^2 - 0.3379\,s + 6.074}
         {s^2 + 59.01\,s + 8.95}
    \label{eq:Gs_ls}
\end{equation}

zatým čo model identifikovaný pomocou optimalizácie parametrov prenosovej funkcie metódou \textit{fminsearch} je

\begin{equation}
    G_{s,\mathrm{fmin}}(s) = 
    \frac{0.0007181\,s^2 - 0.05942\,s + 0.9012}
         {s^2 + 18.94\,s + 1.291}
    \label{eq:Gs_fmin}
\end{equation}

Kvalita zhody medzi modelmi a experimentálnymi dátami bola hodnotená pomocou ukazovateľa RMSE:

\begin{align}
    \mathrm{RMSE}_{\mathrm{LS}} &= 0.036792,
    \label{eq:rmse_ls} \\
    \mathrm{RMSE}_{\mathrm{fmin}} &= 0.026996
    \label{eq:rmse_fmin}
\end{align}

Ako je vidieť, oba modely poskytujú veľmi uspokojivú aproximáciu dynamiky systému v okolí zvoleného pracovného bodu,  
pričom model získaný metódou \textit{fminsearch} dosahuje o niečo menšiu hodnotu RMSE.  
Treba však zdôrazniť, že oba identifikované modely sú len druhého rádu,  
čo predstavuje rozumný kompromis medzi jednoduchosťou a presnosťou pre následný návrh regulácie.

% -------------------------------------------------------------------
% Subsubsection 2.2.5 - Aproximácia identifikovaných modelov
% -------------------------------------------------------------------

\subsubsection{Aproximácia identifikovaných modelov}

Na základe získaných prenosových funkcií možno pozorovať, že ich čitatele obsahujú veľmi malé koeficienty pri prvom a druhom rade.  
Prítomnosť týchto členov je s veľkou pravdepodobnosťou len výsledkom numerických „artefaktov“ vzniknutých pri prevode z diskrétnej prenosovej funkcie na spojitú formu.  
Zahrnutie týchto vyšších rádov v čitateli by výrazne komplikovalo niektoré analytické metódy návrhu regulátorov, napríklad metódu umiestnenia pólov.

Pre účely ďalšej práce sme sa preto rozhodli tieto prenosové funkcie aproximovať jednoduchšou formou.  
Najzrejmejší prístup spočíva v tom, že malé koeficienty pri $s$ a $s^2$ v čitateli zanedbáme a overíme vplyv tejto úpravy pomocou simulácie.

Výsledkom tejto aproximácie sú nasledujúce prenosové funkcie:

\begin{equation}
	G_{s,\mathrm{LS,approx}}(s) =
	\frac{6.074}{s^2 + 59.01s + 8.95}
	\label{eq:Gs_ls_approx}
\end{equation}

póly:
\begin{equation}
	s_{1} = -58.8550, \quad s_{2} = -0.1521
	\label{eq:poles_ls_approx}
\end{equation}

a

\begin{equation}
	G_{s,\mathrm{fmin,approx}}(s) =
	\frac{0.9012}{s^2 + 18.94s + 1.291}
	\label{eq:Gs_fmin_approx}
\end{equation}

póly:
\begin{equation}
	s_{1} = -18.8736, \quad s_{2} = -0.0684
	\label{eq:poles_fmin_approx}
\end{equation}

Z uvedených hodnôt pólov je zrejmé, že obidva modely majú póly s negatívnymi reálnymi časťami, teda ide o stabilné systémy.

Následne bola uskutočnená simulácia odozvy systému analogická tej, ktorá bola uvedená na obrázku~\ref{fig:ts_pb_vent6_spir4_ident_comparison}.  
Porovnanie reálnej odozvy systému TS05 s oboma aproximovanými modelmi je zobrazené na obrázku~\ref{fig:ts_pb_vent6_spir4_ident_approximated_comparison}.

\begin{center}
	\vbox{
		\makebox[\textwidth][c]{%
			\includegraphics[width=1.5\textwidth]{ts_pb_vent6_spir4_ident_approximated_comparison.png}
		}
		\figcaption{Porovnanie skokovej odozvy reálneho systému TS05 a aproximovaných modelov}
		\label{fig:ts_pb_vent6_spir4_ident_approximated_comparison}
	}
\end{center}

Výsledné hodnoty strednej kvadratickej chyby (RMSE) medzi reálnym systémom a simulovanými odozvami boli:

\begin{align}
	\mathrm{RMSE}_{\mathrm{LS,approx}} &= 0.036951,
	\label{eq:rmse_ls_approx} \\
	\mathrm{RMSE}_{\mathrm{fmin,approx}} &= 0.026998
	\label{eq:rmse_fmin_approx}
\end{align}

Z výsledkov, ako aj z priebehov simulácií, je zrejmé, že rozdiel medzi pôvodnými a aproximovanými modelmi je zanedbateľný.  
Hodnota RMSE sa zhoršila len minimálne, čo je z pohľadu ďalšieho použitia modelov úplne prijateľné, najmä vzhľadom na zjednodušenie analytického spracovania, ktoré táto aproximácia prináša.

% -------------------------------------------------------------------
% Subsection - Komunikácia s TS05
% -------------------------------------------------------------------

\subsection{Komunikácia so systémom TS05}
\label{sec:komunikacia_ts05}

V tejto časti je opísaný spôsob komunikácie laboratórneho tepelného systému TS05
s jednotlivými hardvérovými a softvérovými platformami použitými v práci.
Cieľom bolo zabezpečiť jednotný a funkčný prístup k vstupom a výstupom systému
nielen v prostredí \textit{MATLAB/Simulink}, ale aj pri implementácii na platforme
\textit{Arduino UNO} v jazyku C++ a na priemyselnom PLC systéme \textit{Siemens}.

Osobitná pozornosť bola venovaná najmä tomu, aby bolo možné realizovať experimenty
aj mimo pôvodného simulačného prostredia, vytvárať nové riadiace algoritmy
a zároveň spoľahlivo zaznamenávať namerané údaje pre ďalšie spracovanie.

% -------------------------------------------------------------------
% Subsubsection - MATLAB
% -------------------------------------------------------------------

\subsubsection{Komunikácia s MATLAB}
\label{sec:komunikacia_matlab}

Táto časť bude doplnená v nasledujúcej verzii práce.

% -------------------------------------------------------------------
% Subsubsection - C++
% -------------------------------------------------------------------

\subsubsection{Komunikácia v jazyku C++}
\label{sec:komunikacia_cpp}

Keďže systém TS05 bol pôvodne používaný predovšetkým v prostredí
\textit{MATLAB/Simulink} a k dispozícii nebola podrobná technická dokumentácia
popisujúca jednotlivé vodiče a signály, bolo najprv potrebné experimentálne určiť,
ktorý kontakt konektora zodpovedá jednotlivým vstupom a výstupom systému.
Súčasne bolo potrebné vytvoriť nový prepojovací kábel a navrhnúť fyzické zapojenie,
ktoré by umožnilo komunikáciu systému TS05 s platformou \textit{Arduino UNO}.

Zo zariadenia TS05 je vyvedený kruhový 8-pinový konektor.
Pomocou multimetra boli identifikované jednotlivé kontakty, pričom výsledné
priradenie je uvedené v tabuľke~\ref{tab:ts05_pinout}.

\begin{table}[h!]
    \centering
    \caption{Identifikované kontakty konektora systému TS05}
    \label{tab:ts05_pinout}
    \begin{tabular}{|c|l|}
        \hline
        \textbf{Pin} & \textbf{Funkcia} \\ \hline
        1 & Snímač 1 \\ \hline
        2 & Snímač 2 \\ \hline
        6 & Ventilátor \\ \hline
        7 & Špirála \\ \hline
        8 & GND \\ \hline
    \end{tabular}
\end{table}

Po identifikácii pinov bol zhotovený nový kábel, ktorý umožnil priamy prístup
k meraným aj akčným signálom systému TS05. Následne bola zostavená riadiaca
a meracia reťaz založená na mikrokontroléri \textit{Arduino UNO}.

Akčné veličiny systému, teda ventilátor a vykurovacia špirála, boli z platformy
\textit{Arduino} generované pomocou PWM signálu s rozsahom 0 až 255.
Keďže samotný systém TS05 pracuje s analógovými vstupmi v rozsahu 0 až 10~V,
bolo potrebné použiť prevodník PWM/0--10~V.
Tento prevodník vyžaduje externé napájanie minimálne 12~V a 0,1~A.
Takýmto spôsobom bolo možné z digitálnych PWM výstupov \textit{Arduino}
vytvoriť analógové riadiace signály použiteľné pre systém TS05.

Na druhej strane, výstupné signály snímačov systému TS05 sú analógové napätia
v rozsahu 0 až 10~V.
Keďže analógové vstupy platformy \textit{Arduino UNO} pracujú iba v rozsahu
0 až 5~V, bolo potrebné tieto signály vhodne upraviť.
Na tento účel boli použité napäťové deliče, pričom každý delič bol zostavený
z dvoch rezistorov s hodnotou 90~k$\Omega$.
Výstupné napätie deliča je dané vzťahom

\begin{equation}
    U_{\mathrm{out}} = U_{\mathrm{in}} \frac{R_2}{R_1 + R_2}
    \label{eq:voltage_divider}
\end{equation}

pričom pre $R_1 = R_2 = 90~\mathrm{k}\Omega$ dostaneme

\begin{equation}
    U_{\mathrm{out}} = \frac{U_{\mathrm{in}}}{2}
    \label{eq:voltage_divider_half}
\end{equation}

Z toho vyplýva, že signál 0 až 10~V zo snímača je po delení prevedený približne
na rozsah 0 až 5~V, čo už zodpovedá možnostiam analógových vstupov
mikrokontroléra. Samotný prevod v \textit{Arduino} je následne reprezentovaný
hodnotou A/D prevodníka v rozsahu 0 až 1023.

Mapovanie vstupov a výstupov medzi systémom TS05 a platformou \textit{Arduino UNO}
je uvedené v tabuľke~\ref{tab:arduino_mapping_ts05}.

\begin{table}[h!]
    \centering
    \caption{Mapovanie vstupov a výstupov medzi TS05 a Arduino UNO}
    \label{tab:arduino_mapping_ts05}
    \begin{tabular}{|c|l|c|c|}
        \hline
        \textbf{Port Arduino} & \textbf{Signál} & \textbf{Reálny rozsah} & \textbf{Digitálny rozsah} \\ \hline
        A1   & Snímač 1    & 0--10~V & 0--1023 \\ \hline
        A2   & Snímač 2    & 0--10~V & 0--1023 \\ \hline
        PWM5 & Ventilátor  & 0--10~V & 0--255  \\ \hline
        PWM6 & Špirála     & 0--10~V & 0--255  \\ \hline
    \end{tabular}
\end{table}

Celkové fyzické zapojenie systému je znázornené na obrázku~\ref{fig:arduino_ts05_connection}.

\begin{center}
\vbox{
\makebox[\textwidth][c]{%
    
\includegraphics[width=1.2\textwidth]{in_progress.png}
}
\figcaption{Schéma prepojenia systému TS05 s platformou Arduino UNO}
\label{fig:arduino_ts05_connection}
}
\end{center}

Po vytvorení hardvérovej časti bolo potrebné vyriešiť aj interpretáciu vstupných
a výstupných hodnôt.
Najjednoduchším riešením by bolo použiť lineárny prevod medzi digitálnymi rozsahmi
0 až 255, resp. 0 až 1023 a fyzikálnym rozsahom 0 až 10~V.
Experimentálne merania však ukázali, že hoci je správanie systému približne lineárne,
nejde o ideálne lineárny prevod.

Odchýlky sú spôsobené viacerými faktormi.
V prípade napäťových deličov zohráva úlohu tolerancia použitých rezistorov,
ktorá spôsobuje, že reálny deliaci pomer sa mierne odlišuje od ideálnej hodnoty.
Pri PWM/0--10~V prevodníku sa zasa prejavilo, že pri minimálnom použitom napájaní
12~V dochádza k miernemu podhodnocovaniu výstupného napätia.
Z tohto dôvodu boli pre všetky štyri signálové cesty vykonané experimentálne merania,
v ktorých sa porovnávali hodnoty generované alebo čítané mikrokontrolérom
s hodnotami nameranými multimetrom.

Pre každý vstup a výstup bolo realizovaných šesť meraní.
Výsledky pre jednotlivé kanály sú uvedené v tabuľkách~\ref{tab:pwm5_vent},
\ref{tab:pwm6_spirala}, \ref{tab:a1_snimac1} a \ref{tab:a2_snimac2}.

\begin{table}[h!]
    \centering
    \caption{Kalibračné merania pre výstup PWM5 -- ventilátor}
    \label{tab:pwm5_vent}
    \begin{tabular}{|c|c|}
        \hline
        \textbf{PWM5} & \textbf{Namerané napätie [V]} \\ \hline
        0   & 0.00 \\ \hline
        50  & 2.19 \\ \hline
        100 & 4.13 \\ \hline
        150 & 6.03 \\ \hline
        200 & 7.91 \\ \hline
        250 & 9.32 \\ \hline
    \end{tabular}
\end{table}

\begin{table}[h!]
    \centering
    \caption{Kalibračné merania pre výstup PWM6 -- špirála}
    \label{tab:pwm6_spirala}
    \begin{tabular}{|c|c|}
        \hline
        \textbf{PWM6} & \textbf{Namerané napätie [V]} \\ \hline
        0   & 0.00 \\ \hline
        50  & 2.36 \\ \hline
        100 & 4.34 \\ \hline
        150 & 6.24 \\ \hline
        200 & 8.07 \\ \hline
        250 & 9.30 \\ \hline
    \end{tabular}
\end{table}

\begin{table}[h!]
    \centering
    \caption{Kalibračné merania pre vstup A1 -- snímač 1}
    \label{tab:a1_snimac1}
    \begin{tabular}{|c|c|}
        \hline
        \textbf{Hodnota A1} & \textbf{Namerané napätie [V]} \\ \hline
        0   & 0.00 \\ \hline
        365 & 3.68 \\ \hline
        390 & 3.90 \\ \hline
        739 & 7.34 \\ \hline
        904 & 8.95 \\ \hline
        972 & 9.64 \\ \hline
    \end{tabular}
\end{table}

\begin{table}[h!]
    \centering
    \caption{Kalibračné merania pre vstup A2 -- snímač 2}
    \label{tab:a2_snimac2}
    \begin{tabular}{|c|c|}
        \hline
        \textbf{Hodnota A2} & \textbf{Namerané napätie [V]} \\ \hline
        0   & 0.00 \\ \hline
        383 & 3.83 \\ \hline
        406 & 4.06 \\ \hline
        655 & 6.55 \\ \hline
        755 & 7.47 \\ \hline
        769 & 7.68 \\ \hline
    \end{tabular}
\end{table}

Z uvedených hodnôt je zrejmé, že priebehy sú síce blízke lineárnemu správaniu,
avšak nie natoľko, aby bolo možné bezvýhradne použiť jednoduchú lineárnu kalibráciu.
Preto bolo rozhodnuté overiť viacero typov aproximačných funkcií.
Pre každý kanál bolo testovaných päť modelov:
lineárny, kvadratický, kubický, mocninový a exponenciálny.
Kvalita aproximácie bola hodnotená pomocou ukazovateľa RMSE.

Výsledky ukázali, že vo všetkých štyroch prípadoch dosiahla najnižšiu chybu
kubická funkcia.
Z tohto dôvodu bola pre prevod medzi digitálnymi a fyzikálnymi veličinami
zvolená práve kubická aproximácia.
Jej všeobecný tvar možno zapísať ako

\begin{equation}
    y = ax^3 + bx^2 + cx + d
    \label{eq:cubic_calibration}
\end{equation}

kde $x$ predstavuje vstupnú hodnotu a $y$ kalibrovaný výstup.

Pre jednotlivé kanály boli použité nasledujúce kalibračné funkcie:

\begin{equation}
    \mathrm{PWM5}_{\mathrm{vent}}(u) =
    0.05938291u^3 - 0.27988543u^2 + 24.20478979u - 0.52522611
    \label{eq:cal_pwm5}
\end{equation}

\begin{equation}
    \mathrm{PWM6}_{\mathrm{spir}}(u) =
    0.08732906u^3 - 0.39322033u^2 + 22.8496962u - 0.61564759
    \label{eq:cal_pwm6}
\end{equation}

\begin{equation}
    U_{\mathrm{sn1}}(A1) =
    4.64044633 \times 10^{-10} A1^3
    - 8.44254606 \times 10^{-7} A1^2
    + 1.02946107 \times 10^{-2} A1
    + 5.87859135 \times 10^{-5}
    \label{eq:cal_a1}
\end{equation}

\begin{equation}
    U_{\mathrm{sn2}}(A2) =
    -9.09109415 \times 10^{-10} A2^3
    + 8.99002333 \times 10^{-7} A2^2
    + 9.78867803 \times 10^{-3} A2
    - 9.65193864 \times 10^{-5}
    \label{eq:cal_a2}
\end{equation}

Aj keď sa neskôr experimentálne ukázalo, že ani táto aproximácia nie je úplne ideálna,
pre potreby tejto práce bola dosiahnutá presnosť dostatočná.
Jedným z možných spôsobov ďalšieho zlepšenia by bolo vykonanie väčšieho počtu
kalibračných meraní a následná identifikácia presnejšej prevodovej charakteristiky.

Softvérová implementácia na platforme \textit{Arduino UNO} bola realizovaná
v jazyku C++ s využitím rozšírenia \textit{PlatformIO} v prostredí \textit{Visual Studio Code}.
Samotný program bol rozdelený na niekoľko logických častí.
Najprv boli definované hardvérové piny, čas vzorkovania a dĺžka experimentu.
Následne boli implementované pomocné funkcie pre výpočet kubickej kalibrácie
a orezávanie hodnôt do povolených rozsahov.

Dôležitou časťou programu bola realizácia periodického vykonávania riadiacej logiky.
Keďže mikrokontrolér neumožňuje garantovať úplne presné spúšťanie výpočtu
v pevne daných časových okamihoch bez ďalších mechanizmov, bola implementovaná
takzvaná \uv{tick} logika.
Program priebežne sleduje čas od posledného vykonania riadiaceho kroku
a ďalší výpočet sa vykoná až vtedy, keď uplynie požadovaný interval
$T_s = 0{,}1~\mathrm{s}$.
Tým sa zabezpečí, že hlavná riadiaca časť programu beží iba raz za každú periódu
vzorkovania a nie pri každom priechode funkciou \texttt{loop()}.

V každom kroku sa najprv načítajú analógové vstupy A1 a A2,
následne sa pomocou kubických funkcií prevedú na fyzikálne veličiny snímačov.
Po vykonaní samotnej riadiacej logiky sa požadované hodnoty ventilátora a špirály
opäť prevedú pomocou kalibračných funkcií na PWM rozsah 0 až 255,
zaokrúhlia sa a odošlú na výstupné piny mikrokontroléra.

Súčasťou programu boli aj pomocné funkcie,
ktoré slúžia na algebraické spracovanie údajov a zjednodušenie samotnej implementácie.
Konkrétne išlo najmä o funkcie na výpočet kubických aproximačných polynómov
a o funkcie zabezpečujúce obmedzenie vypočítaných hodnôt do povolených rozsahov.
Takéto riešenie umožnilo oddeliť samotnú riadiacu logiku od pomocných výpočtov,
čím sa zvýšila prehľadnosť a zároveň aj jednoduchosť ďalších úprav programu.

Z praktického hľadiska bolo potrebné vyriešiť aj automatický záznam experimentálnych dát.
Priame zapisovanie do súboru z prostredia \textit{Arduino IDE} alebo \textit{PlatformIO}
nie je možné, preto bol použitý sériový port ako komunikačný kanál medzi
mikrokontrolérom a počítačom.
Na strane \textit{Arduino} bola implementovaná jednoduchá stavová komunikácia.
Po zapnutí zariadenia odošle mikrokontrolér cez sériový port správu \texttt{READY}.
Následne externý \textit{Python} skript začne počúvať sériový port a po prijatí
tejto správy odošle príkaz \texttt{START}.
Tým sa spustí samotný experiment aj priebežné odosielanie dát v CSV formáte.
Po uplynutí definovaného času experimentu mikrokontrolér odošle správu \texttt{DONE},
čím sa záznam ukončí a logger uzavrie výstupný \texttt{.csv} súbor.

Dôležitou vlastnosťou implementácie je aj samotná štruktúra súboru \texttt{main.cpp},
ktorá bola navrhnutá tak, aby bolo možné v ďalších experimentoch veľmi jednoducho
meniť iba tú časť programu, ktorá súvisí s konkrétnym riadiacim algoritmom.
V praktickom použití to znamená, že na určenom mieste programu, približne na riadku 102,
sa inicializujú príslušné premenné algoritmu,
zatiaľ čo približne na riadku 154 sa zapisuje samotná riadiaca logika.
Všetky ostatné časti programu, ako napríklad obsluha komunikácie,
časovanie experimentu, prevody veličín a základná štruktúra vykonávania,
môžu zostať nezmenené.
Úpravy sú zvyčajne potrebné iba v prípade,
že sa mení množina premenných odosielaných na logovanie
alebo hlavička zapisovaného CSV súboru.

Týmto spôsobom bola zabezpečená nielen samotná komunikácia medzi systémom TS05
a platformou \textit{Arduino UNO}, ale aj externé spúšťanie experimentov
a automatizované logovanie nameraných údajov.
Takto navrhnutá architektúra sa ukázala ako dostatočne flexibilná
pre implementáciu a testovanie viacerých riadiacich algoritmov.
Úplné zdrojové kódy hlavného programu pre \textit{Arduino UNO}
a externého skriptu pre logovanie dát sú uvedené v prílohách
\ref{lst:main_cpp_appendix} a \ref{lst:logger_py_appendix}.

% -------------------------------------------------------------------
% Subsubsection - Siemens PLC
% -------------------------------------------------------------------

\subsubsection{Komunikácia s PLC Siemens}
\label{sec:komunikacia_siemens}

Táto časť bude doplnená v nasledujúcej verzii práce.